\clearpage\SourceLeaf{22}\ArithmeticHead{15}{R}{INHALT DER ABSCHNITTE I—IV DER D. A.}
\begin{center}\bfseries 6.\end{center}
\noindent\hspace*{\parindent}Den ersten strengen Beweis des Gesetzes gab also \textsc{Gauss}. Er beruht auf der Methode der allgemeinen Induktion und ist dadurch charakterisiert, dass er, ohne aus dem eigentlichsten Gebiete der Theorie der quadratischen Reste herauszutreten, allein mit dem Begriffe eines solchen operiert. Aber er ist recht umständlich, da er acht verschiedene Fälle einzeln zu erledigen nötigt. Mit Verwendung des \textsc{Legendre}schen Symbols gelang es später \textsc{Dirichlet} ihn sehr zu vereinfachen und die acht Fälle auf nur zwei verschiedene zusammenzuziehen. Wenn \textsc{Gauss} einmal\footnote{\textit{Disq. Arithm.} art. 76, W. I, S. 60.} gegen eine Äusserung von \textsc{Waring} einwendet, dass Beweise mehr ex \textls[100]{notionibus} quam ex \textls[100]{notationibus} geschöpft werden müssten,\linebreak eine Maxime, welcher die neuere Mathematik weithin Rechnung getragen hat, so hat er zwar mit seinem ersten Beweise des Fundamentaltheorems dies Ziel aufs beste erreicht, gleichwohl zeigt die so übersichtliche \textsc{Dirichlet}sche Darstellung\begingroup\renewcommand{\thefootnote}{**)}\footnote{\textsc{Crelles} Journal für Mathematik \textbf{47} (1853), S. 139; Werke II, S 121.}\endgroup, wie auch eine geeignete \textls[100]{notatio} von nicht zu unterschätzender Bedeutung sein kann. Spricht sich dies doch deutlich in der Fassung aus, welche zuerst \textsc{Legendre} dem Gesetze in der Gleichung
\[
\left(\frac{p}{q}\right)\cdot\left(\frac{q}{p}\right)=(-1)^{\frac{p-1}{2}\cdot\frac{q-1}{2}}
\]
gab, derzufolge wir es jetzt das Reziprozitätsgesetz der quadratischen Reste nennen und in welcher seine grosse Eleganz erst voll zur Erscheinung kommt. Bei beiden Darstellungen aber war es ein Fall, der eine ganz besondere Schwierigkeit verursachte: die für den Beweis erforderliche Tatsache, dass jede Primzahl \(p>5\) und von der Form \(4n+1\) positiv oder negativ genommen Nichtrest einer kleineren Primzahl sei; diese Tatsache ist für den Fall \(p\equiv5\ (\mathrm{mod.}\ 8)\) sehr einfach nachzuweisen, dagegen ist es nicht so für den andern Fall \(p\equiv1\ (\mathrm{mod.}\ 8)\). Nach \textsc{Gauss}’ eigener Aussage hat er ein volles Jahr mit dieser Schwierigkeit gerungen (W. II, S. 4: per integrum annum me torsit operamque enixissimam effugit), bis es ihm endlich gelang, durch eine einfache und doch tiefe Betrachtung, welche \textsc{Kronecker} mit Recht eine Kraft-\newpage
\SourceLeaf{23}\ArithmeticHead{16}{L}{BACHMANN, ÜBER GAUSS’ ZAHLENTHEORETISCHE ARBEITEN.}
\noindent probe des \textsc{Gauss}schen Genius genannt hat, ihrer am 8. April 1796 (Nr. 2 T.) mächtig zu werden.

Nachdem so \textsc{Gauss} das für zwei Primzahlen \(p,\,q\) geltende Gesetz bewiesen, hat er ihm im art. 133 bereits auch diejenige Allgemeinheit gegeben (am 29. April 1796, Nr. 4 T.), welche mittels des von \textsc{Jacobi}\begingroup\renewcommand{\thefootnote}{*)}\footnotemark\renewcommand{\thefootnote}{***)}\footnotetext{Berliner Monatsberichte 1837, S. 127; Werke VI, S. 254.}\endgroup\ verallgemeinerten \textsc{Legendre}schen Symbols sich ausspricht in der Gleichung
\[
\left(\frac{P}{Q}\right)\cdot\left(\frac{Q}{P}\right)=(-1)^{\frac{P-1}{2}\cdot\frac{Q-1}{2}+\frac{\delta-1}{2}\cdot\frac{\varepsilon-1}{2}},
\]
worin \(P,\,Q\) zwei teilerfremde ungerade Zahlen und \(\delta,\,\varepsilon\) den Vorzeichen von \(P,\,Q\) resp. entsprechend gleich \(\pm1\) sind; wie er denn auch die sogenannten Ergänzungssätze d. i. die obigen Sätze für \(-1\) und 2, die nun aus dem Fundamentaltheoreme aufs neue bewiesen werden konnten (art. 145, nach Nr. 56 T. am 4. Febr. 1797), auf den Fall ausgedehnt hat, dass \(p\) eine zusammengesetzte positive ungerade Zahl ist. So ward es dann auch möglich, den quadratischen Charakter einer beliebigen Zahl in Bezug auf einen beliebigen zu ihr teilerfremden ungeraden Modul zu bestimmen, sowie die Linearformen der sogenannten Teiler und Nichtteiler von \(x^2-a\) d. i. der Primzahlen, von denen die Zahl \(a\) quadratischer Rest bezw. Nichtrest ist, aufzustellen.
